\chapter{Fructosamine}

\section*{What Is This Test?}
Fructosamine measures glycated serum proteins, mainly glycated albumin. It reflects average blood-glucose exposure over approximately the preceding two to three weeks, shorter than the two- to three-month period reflected by hemoglobin A1c.

\section*{Alternative Names}
Glycated Serum Protein; Glycated Albumin-Related Test; Serum Fructosamine.

\section*{Why This Test Is Ordered}
It can help assess recent glycemic change when HbA1c is unreliable because of hemoglobin variants, altered red-cell survival, recent transfusion, pregnancy-related changes, or the need to evaluate a rapid treatment change. It is an adjunct, not a replacement for glucose monitoring in every patient.

\section*{How This Test Is Performed}
A blood sample is collected, usually without special preparation. The laboratory measures glycated proteins by a colorimetric or other validated assay and reports the result in units that are laboratory-specific.

\section*{Preparation, Risks, and Limitations}
Risks are limited to venipuncture. Low albumin, nephrotic protein loss, liver disease, thyroid disease, and major changes in serum-protein turnover can distort the result. Fructosamine is not suitable for interpreting a single episode of hyperglycemia and should be correlated with glucose records.

\section*{Understanding Results}
Higher fructosamine generally indicates higher average glucose during the preceding several weeks. The result is interpreted against the laboratory range and the patient's glucose measurements; it should not be converted to an HbA1c value without acknowledging assay and clinical uncertainty.

\section*{References}
American Diabetes Association guidance on alternative measures of glycemia.

\clearpage
\section*{Understanding Results and Follow-up}
The laboratory report should identify the specimen, method, units, reference interval or decision limit, and whether the result is positive, negative, abnormal, or inconclusive. Reference intervals vary with age, sex, pregnancy, specimen type, assay, and laboratory; a result outside a range is not automatically a diagnosis, and a result within a range does not always exclude disease. Discuss unexpected results with the ordering clinician, who can decide whether repeat or confirmatory testing, treatment, or referral is appropriate. Do not change medicines or delay urgent care based on an isolated result.


\section*{Regulatory Status and Authorization Date}
This chapter describes a test category, not a specific commercial product. FDA, CDSCO, EU IVDR, and MHRA approval or registration dates therefore cannot be assigned without the assay manufacturer, product name, instrument, and intended use. For laboratory-developed tests, an individual agency approval date may not exist. See the Preface for the interpretation of regulatory status.

\clearpage
